\scanpage{18}
cela marchera ($\pi$ est alors \unsure{fidlt|fid.} plat). Or cela suffit ici,
en se ramenant au cas $S=M$ (variété des modules
formels). Donc $H^{10}$ \uline{est stable pour la stratification}.

Donc $H$ admet une décomposition fonctorielle
et stable par ext. de la base
\[
H\simeq H^{10}+H^{01}
\]
et cette décomposition est \uline{stable par stratification}.

Donc, le pb est maintenant de trouver une
stratification \deleted{\unsure{can.}} (fonctorielle ds les deux sens)
sur les deux morceaux $H^{10}=t_A^\vee$ et $H^{01}=t_{A^*}$,
Le deuxième pb, \inserted{pour $A_0$,} se ramène au premier pour $A_0^*$.

Donc il faut trouver une stratification fonctorielle sur
$t_A$. Or il y en a déjà \inserted{une}, grâce à la structure
de puissances p.ièmes non dégénérée, qui
\margin{à préciser \ldots}
% slanted
permet \unsure{de introduire|d'introduire|d'intr. d'emblée} le \unsure{grpe|groupe} structural
discret $\mathrm{Gl}(g)_k$. Et c'est la seule, \deleted{\unsure{(pour $\Lambda,A$ donnés|(pour $A,A$ donnés}} comme on voit
aisément en utilisant la compatibilité avec \unsure{la|les}
puissances p.ièmes (pour $\Lambda$, $A$ sur $\unsure{\lambda|\Lambda|X}$, donnés).

\uline{Remarque}\quad On voit donc que la position de
la filtration p.r. à la stratification ne peut aucunement
donner un principe de classification \inserted{des déf.} infinitésimale en car. $p>0$ !
% archive sheet number: 17
